% ============================================================
%  chapters/ch10-gesture-primary-reality.tex
% ============================================================

\chapter{Gesture as Primary Reality}

\chapterepigraph{%
  The hand is already thinking before the symbol appears.\\
  The trace precedes the glyph.
}{Flyxion, \textit{Gesture Before Symbol}}

\sidenote{See also:\\ Ch.\,12 on keyboards as gesture sensors;\\ Ch.\,28 on gesture-to-glyph}

We have spent three parts building the formal machinery of constraint spaces,
Spherepop collapse, and accessibility landscapes. This chapter turns to the
question that prompted all of it: \emph{where does meaning come from?}

The answer proposed here is gestural. Meaning does not originate in symbols.
Symbols are \emph{compressions} of gestures — stable residues left by repeated
motor acts in constrained environments. The symbol ``A'' is fossilized hand
motion. The word ``grasp'' is a cognitive trace of the act of grasping.

This is not a metaphor. It is, we argue, a claim with formal consequences.

% ── 10.1 ─────────────────────────────────────────────────────
\section{Developmental Evidence}

\begin{historynote}
  McNeill's longitudinal studies \autocite{mcneill1992} established that
  gesture and speech co-originate in development. Children do not learn to
  gesture and then attach words. Gesture and protolanguage co-emerge from the
  same motor substrate. This finding was largely ignored by computational
  linguistics for two decades.
\end{historynote}

In developmental sequence:
\begin{enumerate}
  \item \textbf{Reach-grasp} trajectories appear at 4–6 months — before any
        intentional vocalization.
  \item \textbf{Deictic pointing} emerges at 9–12 months — establishing shared
        attention before shared reference.
  \item \textbf{Iconic gestures} appear at 12–18 months, structurally
        resembling the objects they indicate.
  \item \textbf{Symbolic reference} consolidates at 18–24 months — built
        on a gestural substrate already 12 months old.
\end{enumerate}

The implication is direct: symbolic reference is not the foundation. It is a
\emph{late compression} of a richer gestural system.

\begin{definition}{Gestural Substrate}{gestural-substrate}
  The \emph{gestural substrate} $\gesturesp$ is the space of motor trajectories
  $\gesture{t} : [0, T] \to \mathbb{R}^{3k}$ (for $k$ joints) equipped with:
  \begin{itemize}
    \item an admissibility structure $\admissible{\constraint}$ encoding
          biomechanical constraints,
    \item an equivalence relation $\sim$ identifying motor-equivalent trajectories,
    \item a compression map $\kappa : \gesturesp/{\sim} \to \Sigma^*$ into a
          symbol alphabet $\Sigma$.
  \end{itemize}
  A \emph{symbol} $s \in \Sigma^*$ is any element in the image of $\kappa$.
  The fiber $\fiber{\kappa}{s}$ is the set of gestures that compress to $s$.
\end{definition}

\begin{remark}
  The fiber $\fiber{\kappa}{s}$ is generically non-trivial: many gestures
  produce the same symbol. This is not a defect but a feature — it is precisely
  the robustness of symbolic communication. A child can write ``A'' in many
  different ways, yet the symbol is recognized. The projection $\kappa$ throws
  away all the variation in $\fiber{\kappa}{s}$.
\end{remark}

% ── 10.2 ─────────────────────────────────────────────────────
\section{The Gesture Manifold}

\begin{figure}[h]
  \centering
  \begin{tikzpicture}[scale=1.0]
    % gesture-manifold-3d.tex — fiber projection (layered 2D)
\tikzset{
  fiber/.style={thin, opacity=0.6},
}

  % Top layer: gesture space G
  \draw[RSVPdeep, thick, rounded corners=6pt]
    (-3.5, 0.5) rectangle (3.5, 2.6);
  \node[RSVPdeep, font=\small\scshape] at (-2.8, 2.35) {$\mathcal{G}$};

  % Gesture points — cluster for A (amber)
  \fill[RSVPaccent, opacity=0.85] (-2.8,1.8) circle (3pt);
  \fill[RSVPaccent, opacity=0.85] (-2.0,2.1) circle (3pt);
  \fill[RSVPaccent, opacity=0.85] (-1.4,1.6) circle (3pt);
  % cluster for O (purple)
  \fill[RSVPpurple, opacity=0.85] (-0.5,1.9) circle (3pt);
  \fill[RSVPpurple, opacity=0.85] (0.1,2.2) circle (3pt);
  % cluster for M (teal)
  \fill[RSVPmid, opacity=0.85] (1.2,1.7) circle (3pt);
  \fill[RSVPmid, opacity=0.85] (2.0,1.9) circle (3pt);
  \fill[RSVPmid, opacity=0.85] (2.7,1.6) circle (3pt);

  % Labels
  \node[RSVPaccent, font=\tiny] at (-2.1, 1.05) {motor variants of A};
  \node[RSVPpurple, font=\tiny] at (-0.2, 1.05) {variants of O};
  \node[RSVPmid, font=\tiny] at (1.96, 1.05) {variants of M};

  % Projection arrow
  \draw[-{Stealth[length=7pt]}, RSVPdeep, very thick]
    (0, 0.4) -- (0, -0.5)
    node[midway, right=4pt, font=\small\itshape] {$\kappa$};

  % Bottom layer: symbol space Σ*
  \draw[RSVPdeep, thick, rounded corners=6pt]
    (-3.5,-1.6) rectangle (3.5,-0.7);
  \node[RSVPdeep, font=\small\scshape] at (-2.8,-0.82) {$\Sigma^*$};

  \fill[RSVPaccent] (-2.1,-1.15) circle (4pt);
  \node[RSVPdeep, font=\bfseries\small, below=2pt] at (-2.1,-1.15) {A};
  \fill[RSVPpurple] (-0.2,-1.15) circle (4pt);
  \node[RSVPdeep, font=\bfseries\small, below=2pt] at (-0.2,-1.15) {O};
  \fill[RSVPmid] (1.9,-1.15) circle (4pt);
  \node[RSVPdeep, font=\bfseries\small, below=2pt] at (1.9,-1.15) {M};

  % Fiber lines (hardcoded, no foreach with color vars)
  \draw[fiber, RSVPaccent!70] (-2.8,1.8) -- (-2.1,-0.7);
  \draw[fiber, RSVPaccent!70] (-2.0,2.1) -- (-2.1,-0.7);
  \draw[fiber, RSVPaccent!70] (-1.4,1.6) -- (-2.1,-0.7);
  \draw[fiber, RSVPpurple!70] (-0.5,1.9) -- (-0.2,-0.7);
  \draw[fiber, RSVPpurple!70] (0.1,2.2) -- (-0.2,-0.7);
  \draw[fiber, RSVPmid!70] (1.2,1.7) -- (1.9,-0.7);
  \draw[fiber, RSVPmid!70] (2.0,1.9) -- (1.9,-0.7);
  \draw[fiber, RSVPmid!70] (2.7,1.6) -- (1.9,-0.7);

  % Fiber annotation
  \node[RSVPaccent!80, font=\footnotesize\itshape] at (-3.6, 2.5)
    {$\kappa^{-1}(\texttt{A})$};

  \end{tikzpicture}
  \caption{The gesture manifold $\gesturesp$ with three regions: the
    admissible trajectory space $\Acc(\constraint)$ (outer shell), the
    motor equivalence classes $[\gamma]_\sim$ (middle layer, shown as
    vertical fibers), and the symbolic projection $\Sigma^*$ (the base,
    projected downward). The compression map $\kappa$ collapses each fiber
    to a single symbol.}
  \label{fig:gesture-manifold}
\end{figure}

The gesture manifold has topological structure that symbols cannot capture.
Consider the gesture for the letter ``O'': it can be traced clockwise or
counterclockwise. Both map to the same symbol, but they are topologically
distinct as trajectories — one is a loop oriented $+1$, the other $-1$.

\begin{proposition}{Topological Loss Under Compression}{topo-loss}
  Let $\kappa : \gesturesp/{\sim} \to \Sigma^*$ be a compression map with
  non-trivial fibers. Then $\kappa$ necessarily loses topological invariants
  of order $\geq 1$ (homotopy classes, winding numbers, orientation).
\end{proposition}

\begin{proof}
  Since fibers are non-trivial, distinct homotopy classes of trajectories
  can map to the same symbol. Homotopy is not definable in $\Sigma^*$
  (a discrete set), so the map cannot be homotopy-class-preserving. The
  same argument applies to winding numbers and orientation since $\Sigma^*$
  carries no continuous structure.
\end{proof}

This proposition has a direct pedagogical implication. If we teach children
symbols before gesture, we are teaching them to operate in the projected space
$\Sigma^*$ while withholding information about the richer space $\gesturesp$
from which symbols derive. We are, in effect, teaching the shadow before
teaching the object.

% ── 10.3 ─────────────────────────────────────────────────────
\section{Motor Structure as the Substrate of Meaning}

\begin{conceptaside}{The Compression Hierarchy}
  There is a natural compression sequence running from motor acts to
  cultural artifacts:
  \[
    \underbrace{\text{motor trajectory}}_{\gesturesp}
    \;\xrightarrow{\kappa_1}\;
    \underbrace{\text{gesture class}}_{\gesturesp/{\sim}}
    \;\xrightarrow{\kappa_2}\;
    \underbrace{\text{glyph}}_{\Sigma}
    \;\xrightarrow{\kappa_3}\;
    \underbrace{\text{word}}_{\Sigma^*}
    \;\xrightarrow{\kappa_4}\;
    \underbrace{\text{concept}}_{\mathcal{C}}
  \]
  Each arrow is a lossy projection. Information destroyed at each step is
  not recovered at later steps. The \emph{fiber} at each level contains what
  was discarded.

  \medskip
  The \textbf{Gesture Before Symbol} thesis holds that meaning lives
  primarily in $\gesturesp$ and $\gesturesp/{\sim}$, not in $\Sigma^*$.
  Current NLP systems operate almost exclusively at the $\Sigma^*$ level
  and below — they have, in a precise sense, discarded the meaning while
  keeping only the symbols.
\end{conceptaside}

The implications for artificial intelligence are substantial. A system that
learns only from text has access only to $\kappa_3 \circ \kappa_4$ — the
last two compressions. It sees word co-occurrences but not the motor
equivalence classes that give words their gestural coherence. It knows that
``grasp'' and ``grip'' are related but not \emph{why} — because both compress
gesture classes involving finger flexion under load.

\begin{warning}
  This is not an argument that large language models have no understanding.
  It is an argument that their understanding is \emph{systematically incomplete}
  in a characterizable way: they are missing the sub-$\Sigma^*$ structure of
  $\gesturesp$. The missing structure is not randomly distributed — it concerns
  specifically: spatial reasoning, force and effort, affordance, embodied
  metaphor, and temporal dynamics.
\end{warning}

% ── 10.4 ─────────────────────────────────────────────────────
\section{Evolutionary Evidence}

The priority of gesture over symbol is not only developmental — it appears in
evolutionary sequence as well.

\sidenote{Homo habilis:\\2.6–1.5 Ma}

The hominin hand achieved its current precision grip proportions approximately
2 million years ago with \textit{Homo habilis} — roughly 1.5 million years
before any agreed evidence of symbolic behavior. For most of human evolution,
the hand was the primary cognitive instrument.

\begin{historynote}
  Ibn Rushd (Averroës) argued in \textit{De Anima} that the hand is the
  ``organ of organs'' — the instrument that makes all other instruments
  possible. This anticipates the modern claim that tool use co-evolved with
  language, with the hand as the shared substrate. In the RSVP framework,
  the hand is the primary \emph{accessibility transducer}: it converts
  internal constraint structures into external modifications of the world.
\end{historynote}

The trajectory from:

\[
  \text{grip} \;\to\; \text{tool use} \;\to\; \text{mark-making}
  \;\to\; \text{glyph} \;\to\; \text{alphabet}
\]

is a compression sequence of exactly the form described in the previous
section. Each transition discards variation while preserving a projective
residue.

% ── Exercises ─────────────────────────────────────────────────
\section*{Exercises}

\begin{exercise}
  Let $\Sigma = \{0, 1\}$ and consider the compression map
  $\kappa : \gesturesp \to \Sigma^*$ that maps a trajectory $\gamma$ to
  its Morse encoding (dots and dashes based on duration of contact). What
  information is in the fiber $\fiber{\kappa}{\texttt{SOS}}$? Describe
  three distinct gestures that produce the same Morse sequence.
\end{exercise}

\begin{exercise}[title={Topological}]
  Show that the space of all closed loops in $\mathbb{R}^2$ based at
  the origin is not discrete (unlike $\Sigma^*$). What topological invariant
  distinguishes the clockwise ``O'' from the counterclockwise ``O''?
  Why does this invariant vanish under any map to a discrete set?
\end{exercise}

\begin{exercise}[title={Open problem}]
  Construct a formal model of the compression map $\kappa_1 : \gesturesp \to
  \gesturesp/{\sim}$ using the framework of Riemannian submersions. What
  is the appropriate metric on $\gesturesp$ such that motor-equivalent
  trajectories are at zero distance? Can this metric be learned from
  kinematic data?
\end{exercise}
